\section{Introducción: por qué entrar al mundo de la IA a través de los artículos}

Esta sección establece primero la actitud de estudio para este episodio. Xie Qingchi (谢青池) no presenta los artículos como investigador, sino como responsable de producto, autodidacta de largo plazo y observador de las fronteras de la tecnología, y hace público el recorrido que siguió durante más de un año para leer más de 200 artículos. Esta perspectiva es muy valiosa: no busca convertir cada artículo en una clase de matemáticas, sino preguntar qué frontera cambió cada artículo, qué paradigma abrió, quién lo heredó después y cómo se convirtió en la base de productos y de la industria.

Los 36 artículos de este episodio pueden condensarse en cuatro líneas principales. La primera son los paradigmas de modelos: de GPU, AlexNet, Attention y Transformer a MoE, CoT, LoRA, ReAct y Bitter Lesson. La segunda son la Infra y los datos: ZeRO (Zero Redundancy Optimizer, que reduce la redundancia del entrenamiento al fragmentar los estados del optimizador, los gradientes y los parámetros), Scaling Law, Chinchilla, LAION, RefinedWeb y MegaScale explican por qué los modelos pueden escalarse. La tercera son los modelos de lenguaje: Word2Vec, Google Translate, GPT, BERT, GPT-3, InstructGPT y Tulu 3 explican cómo los LLM se convirtieron en la puerta de entrada actual. La cuarta es la multimodalidad: DeepVideo, las redes de dos flujos, GAN, Diffusion, ViT, CLIP, Stable Diffusion y DiT explican cómo los modelos visuales y generativos confluyeron en la era del Transformer.

\begin{figure}[H]
\centering
\includegraphics[width=0.94\textwidth]{four-part-history-map.png}
\caption{Las cuatro líneas principales de los 36 artículos: los paradigmas de modelos, la Infra/los datos, los modelos de lenguaje y la multimodalidad constituyen en conjunto la historia de la evolución de la IA. Diagrama conceptual propio, elaborado a partir de la descripción del video y de la conversación de 00:19:35--03:56:38.}
\end{figure}

\begin{knowledgebox}{Lectura de la figura: no es una lista de artículos, sino un mapa de fronteras}
Los paradigmas de modelos determinan «qué estructuras pueden aprender»; la Infra y los datos determinan «hasta qué escala pueden crecer»; los modelos de lenguaje determinan «cómo las personas invocan la inteligencia con lenguaje natural»; la multimodalidad determina «cómo los modelos entran al mundo de la visión, el video y la generación». Solo al entrelazarse estas cuatro líneas se formó la industria de la IA actual.
\end{knowledgebox}

\begin{importantbox}{Tesis central de este episodio}
El valor de leer artículos no está en memorizar sus títulos, sino en entrar en contacto directo con el origen de los problemas. Comprender de verdad un artículo clave equivale a saber a qué frontera se enfrentaban los investigadores de su momento, por qué eligieron esa solución, qué parte amplificó después la época y qué parte quedó descartada por el hardware o los datos.
\end{importantbox}

\subsection{Aprender IA con IA: el ciclo de lectura de artículos}

El apartado anterior explicó por qué leer artículos; esta sección aborda cómo sostener esa lectura. Xie Qingchi cuenta que los obstáculos que encontró al principio incluían la base matemática, los artículos en inglés, la polisemia de los términos y la falta de una hoja de ruta. La solución no fue resistir a fuerza de voluntad, sino usar la IA para traducir, explicar, hacer preguntas de seguimiento y organizar, de modo que el artículo dejara de ser un PDF aislado y se convirtiera en un nodo de conocimiento que puede discutirse, repasarse y compartirse.

\begin{figure}[H]
\centering
\includegraphics[width=0.96\textwidth]{paper-reading-loop.png}
\caption{El ciclo de lectura de artículos para aprender IA con IA: la traducción, las preguntas, la hoja de ruta, el repaso y el compartir hacen que el artículo deje de ser solo un PDF. Diagrama conceptual propio, elaborado a partir de la conversación de 00:01:30--00:19:35.}
\end{figure}

\begin{knowledgebox}{Lectura de la figura: la IA no lee por ti, sino que te ayuda a levantar un andamiaje}
La IA puede ayudar a superar las barreras del inglés, de la terminología y de los conocimientos previos, pero no puede sustituir la capacidad de plantear problemas. Una buena forma de leer es encontrar primero el problema, después pedirle a la IA que ayude a traducir, explicar, comparar y cuestionar, y por último volver a situar el artículo en la hoja de ruta histórica.
\end{knowledgebox}

\subsection{Resumen de la sección}

El verdadero tema de EP117 es «cómo entrar al mundo de la tecnología». Los 36 artículos son solo el vehículo; el entrenamiento más profundo consiste en construir un mapa de problemas: saber qué conceptos permanecen estables a largo plazo, qué fronteras están cambiando, qué capacidades provienen del modelo y cuáles provienen de la Infra, los datos y el hardware.
